\documentclass[12pt,a4paper]{article}
\usepackage[utf8]{inputenc}
\usepackage[spanish,es-tabla]{babel}
\usepackage{libertine}
\usepackage{tgheros}
\usepackage{sourcesanspro}
\usepackage[T1]{fontenc}
\usepackage{geometry}
\geometry{margin=2.2cm, headheight=18pt, top=2.5cm, bottom=2.2cm}
\usepackage{xcolor}
\definecolor{primary}{HTML}{1B3A5C}
\definecolor{secondary}{HTML}{2E86AB}
\definecolor{accent}{HTML}{E85D2C}
\definecolor{lightbg}{HTML}{EDF2F7}
\definecolor{tablerow}{HTML}{F0F4F8}
\definecolor{darktext}{HTML}{1A202C}
\definecolor{softrule}{HTML}{CBD5E0}
\usepackage{fancyhdr}
\pagestyle{fancy}
\fancyhf{}
\fancyhead[L]{\sffamily\footnotesize\color{primary!80}MEMORIA DE CALCULO ESTRUCTURAL}
\fancyhead[R]{\sffamily\footnotesize\color{primary!80}Cubierta Metalica -- Sistema Completo}
\fancyfoot[C]{\sffamily\footnotesize\color{primary!60}\thepage}
\renewcommand{\headrulewidth}{0.4pt}
\renewcommand{\headrule}{\hbox to\headwidth{\color{secondary!40}\leaders\hrule height \headrulewidth\hfill}}
\usepackage{titlesec}
\titleformat{\section}[hang]
  {\sffamily\LARGE\bfseries\color{primary}}
  {\thesection\enskip}{0pt}{}
  [\vspace{-0.5em}\color{accent}{\rule{\textwidth}{1.2pt}}\vspace{0.2em}]
\titleformat{\subsection}[hang]
  {\sffamily\Large\bfseries\color{secondary}}
  {\thesubsection\enskip}{0pt}{}
\usepackage{booktabs}
\usepackage{longtable}
\usepackage{colortbl}
\usepackage{array}
\usepackage{amsmath,amssymb}
\newcommand{\tableheadrow}{\rowcolor{primary!15}\sffamily\footnotesize\bfseries\color{primary}}
\newcommand{\tabletotalrow}{\rowcolor{secondary!10}\sffamily\bfseries}
\usepackage{graphicx}
\usepackage{tikz}
\usetikzlibrary{babel}
\usepackage{float}
\usepackage{verbatim}
\usepackage{hyperref}
\hypersetup{colorlinks=true, linkcolor=secondary, urlcolor=accent, citecolor=secondary}
\usepackage{caption}
\captionsetup{font={sf,footnotesize,color=primary!80}, labelfont={bf,color=secondary}, justification=raggedright, singlelinecheck=false, skip=6pt}
\begin{document}

% ═══ PORTADA ═══
\thispagestyle{empty}
\begin{tikzpicture}[remember picture,overlay]
  \fill[primary] (current page.north west) rectangle ([yshift=-2.8cm]current page.north east);
  \fill[secondary!80] (current page.south west) rectangle ([yshift=1.2cm]current page.south east);
  \fill[accent] (current page.north west) rectangle ([yshift=-2.85cm]current page.north east);
\end{tikzpicture}
\vspace*{2cm}
\begin{center}
  {\sffamily\fontsize{16}{20}\selectfont\color{white}\bfseries MEMORIA DE CALCULO ESTRUCTURAL\par}
  \vspace{0.8cm}
  {\sffamily\fontsize{28}{34}\selectfont\color{primary}\bfseries Cubierta Metalica\par}
  {\sffamily\fontsize{22}{28}\selectfont\color{secondary}\bfseries Correas + Cercha de Transicion + Cercha Warren\par}
  \vspace{1.2cm}
  {\sffamily\large\color{darktext!70} Analisis por Elementos Finitos --- NSR-10 LRFD (AISC 360-10)\par}
  \vspace{1.5cm}
  \begin{tikzpicture}[scale=1.1]
    \draw[primary!50, thick] (0,0.8) -- (4,0.8);
    \draw[primary!50, thick] (0,0) -- (4,0);
    \foreach \x in {0,0.5,1,...,4} {
      \draw[secondary!30, thin] (\x,0) -- (\x,0.8);
    }
    \node[primary!40] at (2,-0.3) {\sffamily\scriptsize Cercha Warren 16.52m};
    \draw[accent!50, thick] (0,-0.8) -- (1.5,-0.8);
    \node[accent!40] at (0.75,-1.1) {\sffamily\scriptsize Correa C220};
    \node[secondary!40] at (5.5,-0.4) {\sffamily\scriptsize Cercha transicion 6.40m};
  \end{tikzpicture}
  \vspace{2cm}
  {\sffamily\large\color{white}\bfseries DogCalC v0.4.4 --- Motor PyNite FEM\par}
  {\sffamily\small\color{white!80}\today\par}
\end{center}
\newpage

\pagestyle{fancy}
\renewcommand{\contentsname}{\sffamily\color{primary}\LARGE\bfseries Contenido}
\tableofcontents
\newpage

%============================================
\section{DATOS GENERALES DEL PROYECTO}
%============================================
\begin{center}
\begin{tabular}{>{\sffamily\color{primary!80}}l >{\sffamily}l}
\textbf{Proyecto:}   & Cubierta metalica: correas + cercha transicion + cercha Warren \\
\textbf{Software:}   & DogCalC v0.4.4 (motor PyNite FEM) \\
\textbf{Norma:}      & NSR-10 --- AISC 360-10 para acero \\
\textbf{Metodo:}     & Diseno por Resistencia Ultima (LRFD) \\
\textbf{Acero:}      & Acero Grado 50 ($F_y=345$ MPa, $F_u=450$ MPa) \\
\end{tabular}
\end{center}

La cubierta consta de tres componentes:
\begin{enumerate}
    \item \textbf{Correas} tipo C 220$\times$65$\times$2 mm, L=5.70 m, separacion 1.00 m
    \item \textbf{Cercha de transicion} L=6.40 m (viga de borde), 3 modulos de 6.40 m = 19.20 m total
    \item \textbf{Cercha Warren} en cajon 3D, L=16.52 m, h=1.00 m
\end{enumerate}

Las cerchas de transicion (ejes B y C) transmiten sus reacciones como cargas puntuales
sobre la cercha Warren cajon en los nudos inferiores (N27, N67, N33, N73).
\vspace{0.5cm}

%============================================
\section{PLANTA DE CUBIERTA — DISPOSICION DE EJES}
%============================================
La cubierta se organiza sobre una grilla rectangular de ejes:
\begin{itemize}
    \item \textbf{Ejes 1 a 4} (direccion longitudinal): 3 modulos de 6.40 m = \textbf{19.20 m total}
    \item \textbf{Ejes A a D} (direccion transversal): 16.52 m total (cajon Warren)
\end{itemize}

Dos \textbf{cerchas de transicion} Warren 2D (ejes B y C, L=6.40 m c/u, 3 modulos = 19.20 m)
transmiten las cargas de correas a los \textbf{cajones Warren 3D} (ejes 2 y 3, L=16.52 m).
Las intersecciones B$\times$2, B$\times$3, C$\times$2 y C$\times$3 caen exactamente sobre
nudos existentes de ambas cerchas, sin requerir montantes adicionales.
No existen columnas en el area central (entre B-C y 2-3).

\begin{figure}[H]
\centering
\includegraphics[width=\textwidth]{planta_estructural.png}
\caption{Planta estructural de la cubierta metalica. Ejes 1--4 (19.20 m) y A--D (16.52 m).
Cajones Warren en ejes 2 y 3 (cian). Cerchas de transicion en ejes B y C (naranja).
Las intersecciones $\times$ transfieren las reacciones como cargas puntuales.}
\label{fig:planta-ejes}
\end{figure}

\newpage
% ============================================================
\part*{PARTE 1: CORREA C 220$\times$65$\times$2}
\addcontentsline{toc}{part}{Parte 1: Correa C 220$\times$65$\times$2}
% ============================================================

%============================================
\section{CORREA --- GEOMETRIA}
%============================================
\begin{center}
\begin{tabular}{>{\sffamily\color{primary!80}}l >{\sffamily}l}
\textbf{Perfil:}     & 2C armado 220$\times$65$\times$4 mm (cajon)\\
\textbf{Luz:}        & 5.70 m \\
\textbf{Modelo:}     & Viga simplemente apoyada (3 nodos, 2 elementos) \\
\textbf{Apoyos:}     & N1 PINNED, N2 ROLLER, N3 LATERAL (arriost. $L/2$) \\
\textbf{Ancho aferente:} & 1.00 m \\
\end{tabular}
\end{center}

\subsection{Propiedades de la seccion}
\begin{center}
\begin{tabular}{l>{\sffamily}c}
\toprule
\rowcolor{primary!15}
\multicolumn{1}{c}{\sffamily\footnotesize\bfseries\color{primary} Propiedad} & \multicolumn{1}{c}{\sffamily\footnotesize\bfseries\color{primary} Valor} \\
\midrule
Designacion & 2C 220$\times$65$\times$4 mm (cajon armado) \\
Area total ($A_g$) & 2\,800 mm\textsuperscript{2} \\
Inercia ($I_x$) & $19.68\times10^{6}$ mm\textsuperscript{4} \\
Modulo elastico ($S_x$) & 178.9 cm\textsuperscript{3} \\
Modulo plastico ($Z_x$) & 211.2 cm\textsuperscript{3} \\
Peso propio & 0.216 kN/m \\
\bottomrule
\end{tabular}
\end{center}

\subsection{Modelo estructural}
\begin{figure}[H]
\centering
\begin{tikzpicture}[scale=1.2]
  \draw[thick, primary] (0,0) -- (5.7,0);
  \fill[primary] (0,0) circle (0.08) node[below] {\tiny N1 (PIN)};
  \fill[primary] (2.85,0) circle (0.08) node[below] {\tiny N3};
  \fill[primary] (5.7,0) circle (0.08) node[below] {\tiny N2 (ROLLER)};
  \draw[<->, softrule] (0,-0.4) -- (5.7,-0.4) node[midway,below,font=\tiny\sffamily] {L=5.70 m};
  \node[primary, above, font=\tiny\sffamily] at (1.425,0.15) {M1};
  \node[primary, above, font=\tiny\sffamily] at (4.275,0.15) {M2};
\end{tikzpicture}
\caption{Modelo de la correa. Dos elementos frame con nodo central para lectura de resultados.}
\label{fig:correa-modelo}
\end{figure}

%============================================
\section{CORREA --- CARGAS}
%============================================
Cargas superficiales aplicadas como distribuidas uniformes (kN/m) sobre cada elemento,
con ancho aferente de 1.00 m.

\begin{center}
\begin{tabular}{l c >{\sffamily}c}
\toprule
\tableheadrow
\textbf{Carga} & \textbf{kN/m\textsuperscript{2}} & \textbf{kN/m} \\
\midrule
Peso propio (D)     & 0.80  & 0.80 \\
Cubierta viva (Lr)  & 0.50  & 0.50 \\
\rowcolor{tablerow}
Viento (W)          & 0.40  & 0.40 \\
Granizo (G)         & 1.00  & 1.00 \\
\bottomrule
\end{tabular}
\end{center}

Todos los casos incluyen peso propio estructural (SELFWEIGHT Y -1), que agrega 0.054 kN/m adicionales.

Combinaciones NSR-10 (LRFD): 101 a 108 (identicas a las de la cercha Warren).

%============================================
\section{CORREA --- ANALISIS ESTRUCTURAL}
%============================================
{\sffamily Modelo: 3 nodos, 2 elementos frame (6 GDL/nudo). Apoyos: N1 PINNED, N2 ROLLER.}

\subsection{Reacciones y momentos}
Combo gobernante: \textbf{105} (1.2D+1.6G+0.5W).

\begin{center}
\begin{tabular}{l r r r}
\toprule
\tableheadrow
\textbf{Nudo} & \textbf{FY (kN)} & \textbf{Mz max (kN-m)} & \textbf{DY (mm)} \\
\midrule
N1 (PIN)   & 8.37 & -- & 0.0 \\
\rowcolor{tablerow}
N3 (centro) & -- & 11.93 & $-41.0$ \\
N2 (ROLLER) & 8.37 & -- & 0.0 \\
\midrule
\tabletotalrow
\textbf{Total} & \textbf{16.74} & & \\
\bottomrule
\end{tabular}
\end{center}

Desplazamiento maximo: $\delta_{max} \approx$ \textbf{41.0 mm} $\approx L / 139$ (Combo 105).

\textbf{Momentos maximos por combo:}
\begin{center}
\begin{tabular}{c c}
\toprule
\tableheadrow
\textbf{Combo} & \textbf{$M_u$ (kN-m)} \\
\midrule
101 (1.4D)        &  4.60 \\
102 (1.2D+1.6Lr)  &  7.36 \\
\rowcolor{tablerow}
103 (1.2D+1.6G)   & 10.44 \\
104 (1.2D+1.6Lr+0.5W) & 8.23 \\
\rowcolor{tablerow}
\textbf{105 (1.2D+1.6G+0.5W)} & \textbf{11.93} \\
106 (1.2D+1.0W+0.5Lr) & 6.76 \\
\rowcolor{tablerow}
107 (1.2D+1.0W+0.5G) & 7.72 \\
108 (0.9D+1.0W)    & 4.71 \\
\bottomrule
\end{tabular}
\end{center}

%============================================
\section{CORREA --- CODE CHECK AISC 360-10 / NSR-10}
%============================================
Verificacion de resistencia para la correa armada \textbf{2C 220$\times$65$\times$4 mm}
(dos canales paralelos formando cajon), acero Grado 50 ($F_y=345$ MPa, $F_u=450$ MPa).
Norma NSR-10, titulo F (AISC 360-10). Arriostramiento lateral en $L/2$ (N3: DZ fijo).

\subsection{Propiedades del perfil armado}
\begin{center}
\begin{tabular}{l >{\sffamily}c}
\toprule
\rowcolor{primary!15}
\multicolumn{1}{c}{\sffamily\footnotesize\bfseries\color{primary} Propiedad} & \multicolumn{1}{c}{\sffamily\footnotesize\bfseries\color{primary} Valor} \\
\midrule
Area total ($A_g$)   & 2\,800 mm\textsuperscript{2} \\
Inercia ($I_x$)      & $19.68\times10^{6}$ mm\textsuperscript{4} \\
Modulo elastico ($S_x$) & 178.9 cm\textsuperscript{3} \\
Modulo plastico ($Z_x$) & 211.2 cm\textsuperscript{3} \\
Peso propio          & 0.216 kN/m \\
\bottomrule
\end{tabular}
\end{center}

\subsection{Clasificacion de la seccion (AISC B4.1)}
Para seccion cajon armada (2C), las alas de cada canal individual se verifican como
elementos rigidizados del cajon:
\begin{center}
\begin{tabular}{l c c c c}
\toprule
\rowcolor{primary!15}
\multicolumn{1}{c}{\sffamily\footnotesize\bfseries\color{primary} Elemento} &
\multicolumn{1}{c}{\sffamily\footnotesize\bfseries\color{primary} $b/t$} &
\multicolumn{1}{c}{\sffamily\footnotesize\bfseries\color{primary} $\lambda_p$} &
\multicolumn{1}{c}{\sffamily\footnotesize\bfseries\color{primary} $\lambda_r$} &
\multicolumn{1}{c}{\sffamily\footnotesize\bfseries\color{primary} Clase} \\
\midrule
Ala (cajon) & $65/4=16.25$ & $1.12\sqrt{E/F_y}=27.0$ & $1.40\sqrt{E/F_y}=33.7$ & \textcolor{secondary!80!black}{Compacto} \\
\rowcolor{tablerow}
Alma ($h/t_w$) & $220/4=55.0$ & $3.76\sqrt{E/F_y}=90.5$ & $5.70\sqrt{E/F_y}=137$ & \textcolor{secondary!80!black}{Compacto} \\
\bottomrule
\end{tabular}
\end{center}
$b/t=16.25 < \lambda_p=27.0$ $\rightarrow$ \textbf{Seccion compacta}. Puede desarrollar $M_p$.

\subsection{Capacidad a flexion (AISC F7 -- secciones cajon)}
{\sffamily\small
\textbf{Modulo plastico:} $Z_x = \mathbf{211.2 \times 10^{-6}}$ m\textsuperscript{3} \\
$M_n = M_p = F_y \times Z_x = 345{\times}10^3 \times 211.2{\times}10^{-6} = \mathbf{72.9}$ kN-m \\
$\phi_b M_n = 0.90 \times 72.9 = \mathbf{65.6}$ kN-m \\
\textbf{Ratio:} $14.10 / 65.6 = \mathbf{0.215}$ $\rightarrow$ \textcolor{secondary!80!black}{OK}
}

\subsection{Capacidad a cortante (AISC G2)}
{\sffamily\small
$A_w = 2 \times 220 \times 4 = 1\,760$ mm\textsuperscript{2} \\
$h/t_w = 55 < 1.10\sqrt{k_v E/F_y} = 59.2$ $\rightarrow$ $C_v = 1.0$ \\
$V_n = 0.6 \times 345{\times}10^3 \times 1{,}760{\times}10^{-6} = 364.3$ kN \\
$\phi_v V_n = 0.90 \times 364.3 = \mathbf{327.9}$ kN \\
\textbf{Ratio:} $9.9 / 327.9 = 0.03$ $\rightarrow$ \textcolor{secondary!80!black}{OK}
}

\subsection{Resultados combinados}
\begin{center}
\begin{tabular}{l c c c c}
\toprule
\tableheadrow
\textbf{Miembro} & \textbf{Mu (kN-m)} & \textbf{$\phi$Mn (kN-m)} & \textbf{Ratio} & \textbf{Estado} \\
\midrule
M1 & 14.10 & 65.6 & \textbf{0.215} & \textcolor{secondary!80!black}{OK} \\
\rowcolor{tablerow}
M2 & 14.10 & 65.6 & \textbf{0.215} & \textcolor{secondary!80!black}{OK} \\
\bottomrule
\end{tabular}
\end{center}

\begin{center}
\begin{minipage}{0.85\textwidth}
\color{primary!80}\sffamily\small
\rule{\textwidth}{0.3pt}\\[4pt]
\textbf{Conclusion:} La correa armada 2C 220$\times$65$\times$4 mm \textbf{PASA} holgadamente
el code check AISC 360-10 con un ratio maximo de \textbf{0.215}. La seccion cajon compacta
desarrolla toda su capacidad plastica ($\phi M_n=65.6$ kN-m), muy por encima de la demanda
$M_u=14.10$ kN-m (combo 105). Cortante no critico (ratio 0.03).\\
\rule{\textwidth}{0.3pt}
\end{minipage}
\end{center}

\subsection{Verificacion de deflexiones (servicio)}
{\sffamily\small
Flecha maxima sin mayorar (combo servicio D+Lr):
$w = 0.80+0.216+0.50 = 1.516$ kN/m \\
$\delta = \dfrac{5 w L^4}{384 E I} = \dfrac{5 \times 1.516 \times 5.7^4}{384 \times 200{\times}10^6 \times 19.68{\times}10^{-6}} = \mathbf{8.7}$ mm \\[2pt]
$\delta/L = 8.7/5700 = 1/655 < 1/360$ $\checkmark$ Deflexion aceptable.
}

\newpage
% ============================================================
\part*{PARTE 2: CERCHA DE TRANSICION}
\addcontentsline{toc}{part}{Parte 2: Cercha de Transicion}
% ============================================================
Cercha Warren 2D de 8 paneles que recibe las cargas de las correas (ancho aferente 5.60 m)
y las transmite a la cercha principal. Perfil TUBO 70$\times$70$\times$4 mm en acero Grado 50.

%============================================
\section{TRANSICION --- GEOMETRIA}
%============================================
\begin{center}
\begin{tabular}{>{\sffamily\color{primary!80}}l >{\sffamily}l}
\textbf{Tipo:}      & Cercha Warren 2D, plano Z=0 \\
\textbf{Luz total:} & 6.40 m (modulo unitario) \\
\textbf{Altura:}    & 1.00 m (relacion L/h = 6.4) \\
\textbf{Paneles:}   & 8 paneles de 0.80 m cada uno \\
\textbf{Nodos:}     & 18 (9 cordon superior N1--N9, 9 cordon inferior N10--N18) \\
\textbf{Miembros:}  & 33 (8 CS + 8 CI + 8 diagonales + 9 montantes) \\
\textbf{Apoyos:}    & N10 y N18 PINNED (extremos del cordon inferior) \\
\textbf{Ancho aferente:} & 5.60 m \\
\end{tabular}
\end{center}

\textbf{Nota:} La cercha de transicion real se compone de 3 modulos de 6.40 m cada uno,
totalizando 19.20 m de luz (apoyada en columnas en ejes 1 y 4, y en los cajones Warren
en ejes 2 y 3). El analisis del modulo unitario representa la peor condicion de carga
por ser el caso de apoyo simple mas desfavorable.

\textbf{Patron de diagonales:} Los primeros 4 paneles (x=0 a 3.2 m) tienen diagonales
descendentes (N1$\to$N11, N2$\to$N12, N3$\to$N13, N4$\to$N14). Los ultimos 4 paneles
(x=3.2 a 6.4 m) tienen diagonales ascendentes (N14$\to$N6, N15$\to$N7, N16$\to$N8,
N17$\to$N9), formando una configuraci\'on de V en el centro de la luz en el nodo N14.

\begin{figure}[H]
\centering
\begin{tikzpicture}[scale=0.85, every node/.style={font=\tiny\sffamily}]
  \draw[thick, primary] (0,1.0) -- (6.4,1.0);
  \draw[thick, accent] (0,0) -- (6.4,0);
  \foreach \x in {0,0.8,1.6,2.4,3.2,4.0,4.8,5.6,6.4} {
    \draw[thin, softrule] (\x,0) -- (\x,1.0);
    \fill[primary] (\x,1.0) circle (0.045) node[above] {\tiny \pgfmathprint{int(1+\x/0.8)}};
    \fill[accent] (\x,0) circle (0.045) node[below] {\tiny \pgfmathprint{int(10+\x/0.8)}};
  }
  \foreach \x in {0,0.8,1.6,2.4} {
    \draw[thick, secondary!60] (\x,1.0) -- (\x+0.8,0);
  }
  \foreach \x in {3.2,4.0,4.8,5.6} {
    \draw[thick, secondary!60] (\x,0) -- (\x+0.8,1.0);
  }
  \draw[<->] (0,-0.45) -- (6.4,-0.45) node[midway,below,font=\sffamily] {L = 6.40 m};
  \draw[<->] (6.7,0) -- (6.7,1.0) node[midway,right,font=\sffamily] {h = 1.0 m};
  \node[primary, font=\sffamily\small] at (0.4,1.25) {Cordon Superior (CS)};
  \node[accent, font=\sffamily\small] at (0.4,-0.3) {Cordon Inferior (CI)};
\end{tikzpicture}
\caption{Cercha de transicion Warren 2D. Numeracion de nodos y miembros.}
\label{fig:transicion-geo}
\end{figure}

%============================================
\section{TRANSICION --- CARGAS}
%============================================
Las correas transmiten las cargas de cubierta a los nodos del cordon superior.
El ancho aferente de 5.60 m corresponde a la separacion entre cerchas de transicion.

Carga por nodo interior = $q \times 5.60 \times 0.80$ kN (panel de 0.80 m).
Carga en extremos = mitad del valor interior.

\begin{center}
\begin{tabular}{l c c c}
\toprule
\tableheadrow
\textbf{Carga} & \textbf{Superficial (kN/m\textsuperscript{2})} & \textbf{Nudos interiores (kN)} & \textbf{Extremos N1,N9 (kN)} \\
\midrule
Peso propio (D)  & 0.80 & $-3.584$ & $-1.792$ \\
Cubierta viva (Lr) & 0.50 & $-2.240$ & $-1.120$ \\
\rowcolor{tablerow}
Viento (W) & 0.40 & $-1.792$ & $-0.896$ \\
Granizo (G) & 1.00 & $-4.480$ & $-2.240$ \\
\bottomrule
\end{tabular}
\end{center}

Todos los casos incluyen peso propio estructural (SELFWEIGHT Y -1).
Ocho combinaciones NSR-10 (101--108) identicas a las de la correa y la cercha principal.

%============================================
\section{TRANSICION --- ANALISIS ESTRUCTURAL}
%============================================
Modelo 2D: 18 nudos, 33 elementos frame (6 GDL/nudo), conexiones rigidas.
Apoyos: N10 y N18 PINNED (extremos del cordon inferior).

\subsection{Reacciones por modulo}
\begin{center}
\begin{tabular}{l c c c c c}
\toprule
\tableheadrow
\textbf{Carga} & \textbf{FY N10 (kN)} & \textbf{FY N18 (kN)} & \textbf{Total modulo (kN)} & \textbf{FX N10 (kN)} \\
\midrule
D  & 15.17 & 15.06 & 30.23 & 13.38 \\
Lr &  9.80 &  9.68 & 19.48 &  8.78 \\
\rowcolor{tablerow}
W  &  8.01 &  7.89 & 15.90 &  7.25 \\
G  & 18.76 & 18.64 & 37.40 & 16.45 \\
\bottomrule
\end{tabular}
\end{center}

En el apoyo intermedio (eje 2 o 3), dos modulos adyacentes confluyen:
la reaccion total es la suma de ambos apoyos contiguos
($\approx 2 \times$ el valor de un modulo).

\subsection{Reaccion total en eje 2 (aporte al cajon Warren)}
\begin{center}
\begin{tabular}{l c c c}
\toprule
\tableheadrow
\textbf{Carga} & \textbf{Total en eje 2 (kN)} & \textbf{Por cara del cajon (kN)} & \textbf{Nudos de aplicacion} \\
\midrule
D  & 30.23 & 30.23 & 27, 67 \\
Lr & 19.48 & 19.48 & 27, 67 \\
\rowcolor{tablerow}
W  & 15.90 & 15.90 & 27, 67 \\
G  & 37.40 & 37.40 & 27, 67 \\
\bottomrule
\end{tabular}
\end{center}

Mismo valor en el eje 3 (nudos 33, 73). La componente horizontal FX
se cancela entre modulos adyacentes ($\pm 13$ kN $\rightarrow$ neto $<1$ kN).

\newpage
% ============================================================
\part*{PARTE 3: CERCHA WARREN CAJON 3D}
\addcontentsline{toc}{part}{Parte 3: Cercha Warren Cajon 3D}
% ============================================================

%============================================
\section{WARREN --- GEOMETRIA}
%============================================

\subsection{Vista lateral -- Cara A}
\begin{figure}[H]
\centering
\begin{tikzpicture}[scale=0.5, every node/.style={font=\tiny\sffamily}]
  \draw[thick, primary] (0,1.0) -- (16.52,1.0);
  \draw[thick, accent] (0,0) -- (16.52,0);
  \foreach \x/\xn in {0.00/0.90, 0.90/1.80, 1.80/2.70, 2.70/3.60, 3.60/4.50, 4.50/5.40, 5.40/6.30, 6.30/7.13, 7.13/8.03, 8.03/8.93, 8.93/9.83, 9.83/10.73, 10.73/11.55, 11.55/12.45, 12.45/13.35, 13.35/14.25, 14.25/15.15, 15.15/16.05, 16.05/16.52} {
    \draw[thin, softrule] (\x,0) -- (\x,1.0);
    \draw[thin, secondary!40] (\x,1.0) -- (\xn,0);
  }
  \draw[thin, softrule] (16.52,0) -- (16.52,1.0);
  \draw[<->] (0,-0.4) -- (16.52,-0.4) node[midway,below] {Luz: 16.52 m};
  \draw[<->] (16.52+0.4,0) -- (16.52+0.4,1.0) node[midway,right] {h=1.0 m};
\end{tikzpicture}
\caption{Vista lateral cercha Warren. Cordon superior (azul), inferior (rojo), diagonales (celeste).}
\end{figure}

\subsection{Seccion transversal -- Cajon 3D}
\begin{figure}[H]
\centering
\begin{tikzpicture}[scale=1.4]
  \draw[thick, primary] (0,0) rectangle (0.3,1.0);
  \draw[thick, primary] (2.0,0) rectangle (2.3,1.0);
  \draw[thin, dashed, softrule] (0,0) -- (2.0,0);
  \draw[thin, dashed, softrule] (0,1.0) -- (2.0,1.0);
  \draw[thin, dashed, softrule] (0.3,0) -- (2.3,0);
  \draw[thin, dashed, softrule] (0.3,1.0) -- (2.3,1.0);
  \draw[<->, accent] (0.3,0.5) -- (2.0,0.5) node[midway,above] {\tiny 0.28 m};
  \node at (0.15,-0.25) {\tiny Cara A (z=0)};
  \node at (2.15,-0.25) {\tiny Cara B (z=0.28)};
  \node[right] at (3.0,1.0) {\sffamily\bfseries\color{primary} TUBO 70$\times$70$\times$4};
  \node[right] at (3.0,0.7) {\sffamily\footnotesize $A=1{,}056$ mm$^2$, $I=769{,}472$ mm$^4$};
  \node[right] at (3.0,0.4) {\sffamily\footnotesize $S=21{,}985$ mm$^3$, $r=27.0$ mm};
\end{tikzpicture}
\end{figure}

\subsection{Datos geometricos}
\begin{center}
\begin{tabular}{>{\sffamily\color{primary!80}}l >{\sffamily}l}
\textbf{Tipo:}   & Cercha Warren en cajon 3D (dos caras paralelas) \\
\textbf{Luz:}    & 16.52 m \\
\textbf{Altura:} & 1.00 m \\
\textbf{Separacion caras:} & 0.28 m \\
\textbf{Nodos:}   & 80 (40 por cara) \\
\textbf{Miembros:} & 194 \\
\end{tabular}
\end{center}

%============================================
\section{WARREN --- MATERIALES}
%============================================
Dos tipos de seccion segun el miembro:

\begin{center}
\begin{tabular}{l >{\sffamily}c}
\toprule
\rowcolor{primary!15}
\multicolumn{1}{c}{\sffamily\footnotesize\bfseries\color{primary} Propiedad} & \multicolumn{1}{c}{\sffamily\footnotesize\bfseries\color{primary} Valor} \\
\midrule
\textbf{Cuerdas (CS y CI)} & TUBO 70$\times$70$\times$4 + 2 $\angle$ 3''$\times$1/4'' \\
Area total & 2\,914 mm\textsuperscript{2} \\
Inercia fuerte ($I_x$) & 7\,649$\times$10$^3$ mm\textsuperscript{4} \\
Inercia debil ($I_y$) & 1\,802$\times$10$^3$ mm\textsuperscript{4} \\
Radio de giro min ($r_{min}$) & 24.9 mm \\
\rowcolor{tablerow}
\textbf{Diagonales y montantes} & TUBO 70$\times$70$\times$4 mm \\
Area & 1\,056 mm\textsuperscript{2} \\
Inercia ($I$) & 769\,472 mm\textsuperscript{4} \\
Radio de giro ($r$) & 27.0 mm \\
\midrule
Acero & Grado 50 ($F_y = 345$ MPa, $F_u = 450$ MPa) \\
\bottomrule
\end{tabular}
\end{center}

La seccion compuesta de las cuerdas consiste en el tubo cuadrado 70$\times$70$\times$4 mm
con dos angulos L 3''$\times$1/4'' (76$\times$76$\times$6.35 mm) soldados a las caras
opuestas, formando un perfil asimetrico con la mayor inercia en el plano de la cercha.

%============================================
\section{WARREN --- CARGAS Y COMBINACIONES}
%============================================
Cargas aplicadas en nudos del cordon superior de ambas caras,
ancho aferente de \textbf{0.65 m} a cada lado del eje (1.30 m total).

\begin{center}
\begin{tabular}{l c >{\sffamily}c}
\toprule
\tableheadrow
\textbf{Carga} & \textbf{kN/m\textsuperscript{2}} & \textbf{Total cubierta (kN)} \\
\midrule
Peso propio (D)  & 0.80 & 34.36 \\
Cubierta viva (Lr) & 0.50 & 21.48 \\
\rowcolor{tablerow}
Viento (W) & 0.40 & 17.18 \\
Granizo (G) & 1.00 & 42.95 \\
\bottomrule
\end{tabular}
\end{center}

Adicionalmente, las cerchas de transicion (ejes B y C) transmiten sus reacciones
al cajon en los nudos inferiores de cada cara:

\begin{center}
\begin{tabular}{l c c c}
\toprule
\tableheadrow
\textbf{Carga} & \textbf{Reaccion por interseccion (kN)} & \textbf{Por cara del cajon (kN)} & \textbf{Nudos} \\
\midrule
D  & 30.23 & 30.23 & 27, 67, 33, 73 \\
Lr & 19.48 & 19.48 & 27, 67, 33, 73 \\
\rowcolor{tablerow}
W  & 15.90 & 15.90 & 27, 67, 33, 73 \\
G  & 37.40 & 37.40 & 27, 67, 33, 73 \\
\bottomrule
\end{tabular}
\end{center}

Las cargas de transicion se aplican como \textbf{JOINT LOAD FY} en los nudos del cordon inferior
de cada cara (N27 y N67 en interseccion con eje B, N33 y N73 con eje C).
Ocho combinaciones NSR-10 (101--108) identicas a las de la correa.

%============================================
\section{WARREN --- ANALISIS}
%============================================
Modelo 3D: 80 nudos, 194 elementos frame (6 GDL/nudo), conexiones rigidas.
Apoyos: N21 y N61 PINNED, N40 y N80 ROLLER.

\subsection{Reacciones -- Combinacion critica 105 (1.2D+1.6G+0.5W)}
\begin{center}
\begin{tabular}{l r r r}
\toprule
\tableheadrow
\textbf{Nudo} & \textbf{FY (kN)} & \textbf{Tipo} \\
\midrule
N21 (PINNED)  & 143.75 & Cara A, izquierdo \\
N61 (PINNED)  & 143.75 & Cara B, izquierdo \\
N40 (ROLLER)  & 138.46 & Cara A, derecho \\
\rowcolor{tablerow}
N80 (ROLLER)  & 138.46 & Cara B, derecho \\
\midrule
\tabletotalrow
\textbf{Total} & \textbf{564.42} & \\
\bottomrule
\end{tabular}
\end{center}

La carga total de 564.42 kN se distribuye simetricamente entre ambas caras del cajon
y entre apoyos PINNED (57.5 kN diferencia por el empuje horizontal de las diagonales).

\subsection{Reacciones por carga individual}
\begin{center}
\begin{tabular}{l r r r r r}
\toprule
\tableheadrow
\textbf{Carga} & \textbf{N21} & \textbf{N61} & \textbf{N40} & \textbf{N80} & \textbf{Total} \\
\midrule
D  & 41.83 & 41.83 & 40.29 & 40.29 & 164.25 \\
Lr & 27.60 & 27.60 & 26.58 & 26.58 & 108.37 \\
\rowcolor{tablerow}
W  & 22.86 & 22.86 & 22.01 & 22.01 &  89.75 \\
G  & 51.32 & 51.32 & 49.44 & 49.44 & 201.52 \\
\bottomrule
\end{tabular}
\end{center}

%============================================
\section{WARREN --- FUERZAS AXIALES}
%============================================
Envolvente de diseno (combos 101--108) con cargas de cubierta + transicion.
$(+)$ traccion, $(-)$ compresion.

\subsection{Cordon Superior (compresion)}
\begin{center}
\begin{tabular}{r >{\sffamily}r | r >{\sffamily}r}
\toprule
\tableheadrow
\textbf{Miembro} & \textbf{Axial (kN)} & \textbf{Miembro} & \textbf{Axial (kN)} \\
\midrule
M1 & $-146.6$ & M11 & $-730.0$ \\
M2 & $-278.4$ & M12 & $-704.6$ \\
M3 & $-393.3$ & M13 & $-722.2$ \\
M4 & $-490.7$ & M14 & $-725.8$ \\
M5 & $-570.4$ & M15 & $-718.4$ \\
M6 & $-632.5$ & M16 & $-731.9$ \\
M7 & $-677.0$ & M17 & $-731.9$ \\
M8 & $-703.6$ & M18 & $-730.8$ \\
M9 & $-713.0$ & M19 & $-730.0$ \\
M10 & $-704.6$ & & \\
\bottomrule
\end{tabular}
\end{center}
\textbf{Max compresion:} M16/M17 = \textcolor{accent}{$-731.9$ kN} (en las cercanias del
apoyo). Se observa que las cargas de transicion aumentan significativamente las fuerzas
en las cuerdas cercanas a los apoyos.

\subsection{Cordon Inferior (traccion)}
\begin{center}
\begin{tabular}{r >{\sffamily}r | r >{\sffamily}r}
\toprule
\tableheadrow
\textbf{Miembro} & \textbf{Axial (kN)} & \textbf{Miembro} & \textbf{Axial (kN)} \\
\midrule
M20 & $+22.3$ & M30 & $+723.3$ \\
M21 & $+128.0$ & M31 & $+730.2$ \\
M22 & $+263.0$ & M32 & $+729.9$ \\
M23 & $+381.0$ & M33 & $+724.8$ \\
M24 & $+482.0$ & M34 & $+716.8$ \\
M25 & $+565.0$ & M35 & $+706.0$ \\
M26 & $+630.0$ & M36 & $+691.0$ \\
M27 & $+678.0$ & M37 & $+667.0$ \\
M28 & $+706.0$ & M38 & $+619.0$ \\
M29 & $+716.0$ & & \\
\bottomrule
\end{tabular}
\end{center}
\textbf{Max traccion:} M31 = \textcolor{secondary}{$+730.2$ kN}.

\subsection{Diagonales}
Las diagonales trabajan en traccion/compresion con rangos de $\pm 45$ a $\pm 90$ kN,
similar a la configuracion original sin cargas de transicion. Las diagonales cercanas
a los nuevos puntos de carga (nudos 27, 33, 67, 73) presentan un ligero incremento
($+5$ a $+10$ kN).

%============================================
\section{WARREN --- CODE CHECK AISC 360-10 / NSR-10}
%============================================
Verificacion de los 194 miembros de la cercha segun AISC 360-10 (NSR-10 Titulo F).
Acero Grado 50 ($F_y=345$ MPa, $F_u=450$ MPa). Perfil TUBO 70$\times$70$\times$4 mm
($A=1{,}056$ mm\textsuperscript{2}, $r=27.0$ mm, $I=769{,}472$ mm\textsuperscript{4}).

\subsection{Capacidad a traccion (AISC D2)}
{\sffamily\small
\textbf{Fluencia (D2-1):}
$\phi_t P_n = 0.90 \times F_y \times A_g = 0.90 \times 345{\times}10^3 \times 1{,}056{\times}10^{-6} = \mathbf{328}$ kN
\\[2pt]
\textbf{Rotura (D2-2):}
$\phi_t P_n = 0.75 \times F_u \times A_e \approx 0.75 \times 450{\times}10^3 \times 1{,}056{\times}10^{-6} = \mathbf{356}$ kN
\\[2pt]
Gobierna fluencia: $\phi_t P_n = \mathbf{328}$ kN.
}

\subsection{Capacidad a compresion (AISC E3)}
{\sffamily\small
Cordon superior ($L=0.90$ m, $K=1.0$):
\\[2pt]
$KL/r = 1.0 \times 0.90 / 0.0270 = 33.3$ \\
$F_e = \pi^2 E / (KL/r)^2 = \pi^2 \times 200{,}000 / 33.3^2 = \mathbf{1{,}785}$ MPa \\
$F_y / F_e = 345 / 1{,}785 = 0.193$ \\
$F_{cr} = \left[0.658^{F_y/F_e}\right] F_y = 0.920 \times 345 = \mathbf{317}$ MPa \\
$\phi_c P_n = 0.90 \times F_{cr} \times A_g = 0.90 \times 317{\times}10^3 \times 1{,}056{\times}10^{-6} = \mathbf{302}$ kN
}

\subsection{Resultados -- miembros criticos}
\begin{center}
\begin{tabular}{l l c c c c c}
\toprule
\tableheadrow
\textbf{Miembro} & \textbf{Seccion} & \textbf{Tipo} & \textbf{L (m)} & \textbf{Pu (kN)} & \textbf{$\phi$Pn (kN)} & \textbf{Ratio} \\
\midrule
M87  & TUBO+2L & CS-comp & 0.90 & $-758.1$ & 822 & \textbf{0.922} \\
M10  & TUBO+2L & CS-comp & 0.90 & $-758.1$ & 822 & 0.922 \\
M11  & TUBO+2L & CS-comp & 0.90 & $-757.9$ & 822 & 0.922 \\
M88  & TUBO+2L & CS-comp & 0.90 & $-757.9$ & 822 & 0.922 \\
\rowcolor{tablerow}
M29  & TUBO+2L & CI-trac & 0.90 & $+756.8$ & 905 & \textbf{0.836} \\
M106 & TUBO+2L & CI-trac & 0.90 & $+756.8$ & 905 & 0.836 \\
M30  & TUBO+2L & CI-trac & 0.90 & $+756.2$ & 905 & 0.836 \\
M107 & TUBO+2L & CI-trac & 0.90 & $+756.2$ & 905 & 0.836 \\
\bottomrule
\end{tabular}
\end{center}

\subsection{Resumen general}
\begin{center}
\begin{tabular}{l r}
\toprule
\tableheadrow
\textbf{Concepto} & \textbf{Valor} \\
\midrule
Total miembros & 194 \\
OK & \textcolor{secondary!80!black}{194} \\
\rowcolor{tablerow}
SOBREESFORZADOS & \textcolor{accent}{\textbf{0}} \\
Peor compresion & \textbf{0.922} (M87/M10, CS) \\
\rowcolor{tablerow}
Peor traccion & \textbf{0.836} (M29/M106, CI) \\
\bottomrule
\end{tabular}
\end{center}

{\sffamily\small\color{secondary!80!black}
\textbf{Nota:} Las cuerdas fueron reforzadas con 2 angulos L 3''$\times$1/4''
soldados al tubo 70$\times$70$\times$4 mm. La seccion compuesta resultante
($A=2{,}914$ mm$^2$) aumenta la capacidad axial de 302 kN a 822 kN en compresion
y de 328 kN a 905 kN en traccion, suficiente para las cargas de las cerchas
de transicion. Las diagonales y montantes (TUBO 70$\times$70$\times$4 sin reforzar)
permanecen dentro de capacidad con ratios menores a 0.50.
}

%============================================
\section{CONCLUSIONES GENERALES}
%============================================
\begin{enumerate}
    \item \textbf{Correa armada 2C 220$\times$65$\times$4 mm:} \textcolor{secondary!80!black}{PASA} holgadamente.
          Seccion cajon compacta, $\phi M_n=65.6$ kN-m, demanda $M_u=14.10$ kN-m (combo 105),
          ratio \textbf{0.215}. Arriostramiento lateral en $L/2$. Deflexion $L/655$ aceptable.
    \item \textbf{Cercha de transicion Warren 2D:} Todos los 33 miembros OK con TUBO 70$\times$70$\times$4 mm
          Grado 50. Ratio maximo \textbf{0.267} (cordon superior a compresion). Reaccion por modulo
          de 30.23 kN D, 19.48 kN Lr, 15.90 kN W y 37.40 kN G, transmitidas a la cercha principal.
    \item \textbf{Cercha Warren 3D con cuerdas reforzadas:} \textcolor{secondary!80!black}{PASA} con la
          seccion compuesta TUBO 70$\times$70$\times$4 + 2$\angle$3''$\times$1/4''.
          Las cargas de transicion incrementaron las fuerzas axiales a $\sim$760 kN,
          pero la seccion reforzada ($A=2{,}914$ mm$^2$) proporciona capacidad de
          822 kN en compresion y 905 kN en traccion. \textbf{194/194 miembros OK},
          peor ratio \textbf{0.922} (compresion en cuerda superior).
    \item \textbf{Rediseno completado:} La adicion de dos angulos L 3''$\times$1/4'' soldados
          a las caras opuestas del tubo 70$\times$70$\times$4 en las cuerdas del cajon Warren
          resuelve el incremento de carga inducido por las cerchas de transicion.
          Las diagonales y montantes (tubo simple) no requieren refuerzo.
\end{enumerate}

% ═══════════════════════════════════════════
\appendix
\section{Archivos de entrada STAAD (.std)}
% ═══════════════════════════════════════════

\subsection{Correa C 220$\times$65$\times$2}
{\scriptsize\ttfamily\verbatiminput{../examples/correa_c220.std}}

\subsection{Cercha de Transicion Warren 8P}
{\scriptsize\ttfamily\verbatiminput{../examples/cercha_transicion.std}}

\subsection{Cercha Warren Cajon 3D (con cargas de transicion)}
{\scriptsize\ttfamily\verbatiminput{../examples/cercha_warren.std}}

\end{document}
